\section{Threats to Validity}
\label{sec:threats}

\paragraph{Construct validity.}
The evaluation is a \emph{controlled synthetic simulation}, not a
deployment against live LLM APIs. We do not claim that the absolute
utility numbers reflect what a real LLM agent would achieve on
ToolBench or $\tau$-bench. What we claim is that under the simulated
cost/latency distributions, the ordering of policies is consistent,
and the qualitative safety gap (order of magnitude in overrun rate) is
robust across the sweep.

\paragraph{Internal validity.}
Confidence parameters $\alpha = 1.0$, $\beta = 1.5$, $\gamma = 1.5$
are fixed across all cells and were chosen before the final sweep; we
do not report post-hoc tuning. A one-at-a-time sensitivity sweep over
all three parameters (Appendix~\ref{app:abl-abg}) confirms that the
budget-overrun rate remains identically zero across the grid. The only
intentional sensitivity is the $z$ sweep over \CBUC~variants, which is
itself the object of study.

\paragraph{External validity.}
Cost distributions are Pareto-parameterised; real LLM API pricing is
step-functioned by token count and partially deterministic per-call.
Latency is Gaussian; real-world latency is often heavy-tailed with
occasional timeouts. Our calibrated benchmark replay
(Appendix~\ref{app:abl-replay}), parameterised to published summary
statistics of ToolBench, $\tau$-bench, and BFCL, shows that the
qualitative safety gap persists under these workload profiles;
however, a full replay harness over raw API traces remains future work
and is explicitly flagged as \emph{requires empirical validation}
before any production claim.

\paragraph{Chance-constraint approximation.}
The sub-Gaussian correction in \S\ref{sec:algorithm} is conservative
for heavy-tailed cost; under Pareto tails with $\alpha \le 2$ the
variance is infinite, so the theorem's sub-Gaussian proxy is formally
inapplicable. The $\alpha=1.5$ rows in Table~\ref{tab:tail} are
therefore an empirical robustness stress test, not covered theory.
\CBUC's low overrun in those rows appears to come from the combination
of chance-constraint stopping and conservative feasibility checks, not
from a valid Gaussian tail model. Replacing the Gaussian $z$ with a
Chernoff-style, median-of-means, or robust VaR/CVaR correction is
required before claiming theoretical safety under heavy tails.

\paragraph{Oracle definition.}
Our regret is measured against the informed \emph{non-adaptive}
oracle. The adaptivity-gap remark (Remark~\ref{rem:gap}) applies only
under additive-utility stochastic-knapsack conditions; outside those
conditions, adaptive-history dependence remains an open modelling gap.

\paragraph{Statistical significance.}
Table~\ref{tab:headline} reports approximate $95\%$ confidence
intervals, which are sufficient for the large safety gaps but still
informal for utility comparisons. A paired-seed Wilcoxon signed-rank
test across workload cells should be added before camera-ready,
especially for the smaller utility gaps among pre-check UCB,
CVaR-\BwK, \BwK-vanilla, and \CBUC~variants. The ablation studies
(Appendix~\ref{app:ablations}) further corroborate these findings
across $1{,}920$ additional replicate rows.

\paragraph{Baseline strength.}
Concerns that \ReAct-capped alone is an insufficient comparator are
legitimate for an algorithmic-efficacy claim. We therefore include
four constraint-aware baselines drawn from standard families: a
Lagrangian primal-dual safe-CMDP-style controller, a CVaR-aware \BwK,
a deterministic pre-check UCB guarded greedy, and a non-adaptive
chance-constrained knapsack oracle
(Table~\ref{tab:headline}). \CBUC~remains the only policy in this
extended set to keep both constraint-violation rates below $2\%$
uniformly across the sweep. The remaining utility gap to the
Lagrangian baseline is reported as a real deployment trade-off rather
than hidden by the safety aggregate.
